% TOP_PROBLEM: {"id": 3325, "title": "Universal point sets for planar graphs", "category_id": 3, "category": {"id": 3, "name": "graph_theory", "display_name": "Graph Theory", "description": "Problems involving graphs, networks, and their properties.", "slug": "graph-theory", "order_index": 3, "created_at": "2026-07-31T15:26:25.671Z"}, "status": "open", "problem_number": "OPG-326", "set_id": 12}
% FIRST_POSTED: 2026-10-02
% ATTEMPT_STATUS: unresolved
% MODEL: GPT 6 Astra Ultra
% UnsolvedMath ID 3325; source catalogue code OPG-326
% !TeX program = lualatex
% Research attempt dated 2026-10-02; canonical statement preserved.
\documentclass[11pt,a4paper]{article}
\usepackage[margin=25mm]{geometry}
\usepackage{fontspec}
\setmainfont{lmroman10-regular.otf}[BoldFont=lmroman10-bold.otf,
 ItalicFont=lmroman10-italic.otf,BoldItalicFont=lmroman10-bolditalic.otf]
\usepackage{amsmath,amssymb,mathtools}
\usepackage{unicode-math}
\setmathfont{latinmodern-math.otf}
\AtBeginDocument{\let\setminus\smallsetminus}
\usepackage{xurl,hyperref}
\hypersetup{colorlinks=true,urlcolor=blue}
\setcounter{secnumdepth}{0}
\setlength{\parindent}{0pt}
\setlength{\parskip}{5pt}
\setlength{\emergencystretch}{3em}
\begin{document}
\section*{Universal point sets for planar graphs}\label{3325}
{\small UnsolvedMath ID 3325 · OPG-326 · Graph Theory · OpenGarden}
\subsection{Definitions and mathematical statement}
Fix an integer $n\ge1$. A finite point set $P\subset\mathbb R^2$ is
$n$-universal for planar graphs if every finite simple planar graph
$G$ with exactly $n$ vertices has an injective map
$\varphi:V(G)\to P$ such that drawing each edge $uv$ as the straight
segment $[\varphi(u),\varphi(v)]$ gives a plane embedding. Thus edge
interiors are disjoint, no edge contains a vertex other than its
ends, and adjacent edges meet only at their common endpoint.

The graph can use any $n$ points of $P$, and the assignment of
vertices to those points is chosen separately for each graph.
The points themselves must be fixed independently of $G$.
Define $u(n)$ as the smallest possible cardinality of such a point
set. The question asks whether
\[
                  \exists C>0\ \forall n\ge1:\qquad u(n)\le Cn.
\]
Equivalently, is there a family of universal point sets whose sizes
grow at most linearly with $n$? The constant is absolute and no
particular grid, coordinate precision, or algorithm for constructing
the sets is required by the statement.

The graph comes with no prescribed embedding, so its rotation
system may be chosen along with the drawing. The straight-segment
requirement is essential: arbitrary curved edges would allow many
more placements. Allowing unused points means this is not the
stronger demand for a universal set of exactly $n$ points, and no
fixed labelling of points by graph vertices is assumed.
\subsection{Short English statement}
Is there a point set of at most Cn points that supports a crossing-free straight-line drawing of every n-vertex planar graph?
\subsection{Research attempt}
\paragraph{Scope and method.}
This research attempt by GPT-6 Astra Ultra is dated 2 October 2026.
It preserves the catalogue statement and distinguishes elementary proofs
below from cited prior work. No novelty claim or formal proof-assistant
verification is made. The final paragraph states the precise remaining scope.
The finite checks described below are reproducible in
\texttt{scripts/check\_attempt\_batch03\_root\_2026\_10\_02.py}.
\paragraph{The linear target and a recent result.}
The point set must be chosen before the planar graph is known.
Giving each graph its own set of $n$ points does not address
that quantifier order. Gordon's September 2026 preprint [R1]
constructs universal sets of size $n^{1+o(1)}$. This is a major
improvement over the previously quadratic general upper bound,
but an $n^{1+o(1)}$ estimate alone does not supply a fixed
constant $C$ with size at most $Cn$. For example,
$n\exp(\sqrt{\log n})$ has the former form and not the latter.
This catalogue entry therefore asks for a stronger bound than
the subquadratic target in the earlier universal-point-set
entry. We do not reproduce or claim the preprint's theorem.

\paragraph{A sharp positive case: all trees.}
Fix any $n$ points in strictly convex position. We prove that
they simultaneously support every tree on $n$ vertices.
Root a tree and order the children of each vertex arbitrarily.
List the vertices in depth-first preorder, visiting a vertex
before all its children and completing a child's entire
subtree before moving to the next child. Place this order
around the convex polygon and draw tree edges as segments.
The vertices in each rooted subtree occupy a consecutive
interval. Two chords of a strictly convex polygon cross in
their interiors exactly when their four endpoints alternate
in the cyclic order. That cannot happen for tree edges in
the preorder placement. Indeed the subtrees of different
children occupy disjoint intervals. An edge from a parent
to a child's root has every earlier child's subtree between
its endpoints in the linear order; edges internal to those
subtrees have both endpoints there. Edges inside the selected
child's subtree have both endpoints at or after its root,
and edges in later child subtrees lie still later. Applying
the same reasoning recursively treats every edge pair.
No edge passes through another vertex, since strict convexity
puts no three of the chosen points on one line. This proves
the desired simultaneous embedding for the tree class.
At least $n$ points are needed for an injective placement,
so the size is optimal within that class.

\paragraph{Why more points on a convex curve do not suffice.}
Every subset of a strictly convex point set is strictly
convex. Four points chosen from it have two crossing
diagonals, so they cannot support a crossing-free straight
drawing of $K_4$. For every $n\ge4$, the planar graph formed
by $K_4$ and $n-4$ isolated vertices is therefore an
obstruction to universality for any strictly convex point
set, regardless of its cardinality. The tree construction
fails because general planar graphs need interior vertices,
not because the convex curve contains too few candidates.
This gives a geometric obstruction to an entire class of
linear-size proposals without being a lower bound against
all possible point configurations.

\paragraph{A small exact general case.}
For $n=4$, use three corners of a triangle and one strictly
interior point. Its six straight segments give a planar
drawing of $K_4$. Every graph on four vertices is a subgraph
of that complete graph, so this single set is universal
and the optimum size is four. For $n\le3$ a convex set of
$n$ points works as well. Beyond this, combining separate
embeddings by taking their union is not a linear bound:
the number of graphs grows, and there is no argument here
forcing their vertex locations to be shared.

\paragraph{Verification and final status.}
The root script enumerates all labelled trees through six
vertices via Pruefer sequences, constructs the preorder
placement on the integer parabola $(i,i^2)$, and tests all
independent edge pairs with exact determinants. It also
checks the four-point interior construction and the convex
$K_4$ obstruction. The proof above establishes the tree
case for all $n$; enumeration is only an implementation
check. A construction of one set with a uniform constant
times $n$ points for all planar graphs remains absent.
The linear-size question is unresolved by this attempt.

\subsection{Sources}
\begin{itemize}
\item[S1] ULAM.AI and the UnsolvedMath Contributors, supplied problems.json, record 3325 (OPG-326), OpenGarden collection. Catalogue identity and source transcription; CC BY 4.0.
\url{https://huggingface.co/datasets/ulamai/UnsolvedMath}.
\item[S2] Open Problem Garden, Universal point sets for planar graphs. Original problem and discussion; the mathematical formulation below is expanded from the supplied source record.
\url{http://www.openproblemgarden.org/op/small_universal_point_sets_for_planar_graphs}.

\item[R1] T. Gordon, Almost Linear Universal Point Sets for Planar Graphs,
arXiv:2609.10916, submitted 10 September 2026.
\url{https://arxiv.org/abs/2609.10916}.

\end{itemize}
\end{document}
